\documentclass[a4paper, 11pt, UTF8]{ctexart}%{article} % For LaTeX2e
\usepackage{nips13submit_e,times}
\usepackage{hyperref}
\usepackage{url}
\usepackage{graphicx}
\usepackage{color}
\usepackage{multirow}
\usepackage{xcolor}
\usepackage{ulem}
\usepackage{tikz}
\usepackage{eso-pic}
\usepackage{colortbl}
\usepackage{supertabular}
\usepackage[section]{placeins}
\usepackage{float}
\definecolor{t_gary}{gray}{.9}
\title{
\begin{Large}
Face Search Benchmark: internal release vs public solutions
\end{Large}\\
人脸搜索基准测试：内部方案vs公开方案\\
\begin{Large}
  {\color{red} Internal Reference Only 仅供内部参阅}
\end{Large}}
\author{
张继\thanks{最后修改于\today。} \\
Glasssix Research\\
成都，610051 \\
\texttt{\href{mailto:zj@glasssix.com}{zj@glasssix.com}} \\
\And
何贤昆\thanks{最后修改于--。} \\
Glasssix Engineering \\
成都，610051 \\
\texttt{\href{mailto:hermon@glasssix.com}{hermon@glasssix.com}} \\
%\AND
%Coauthor \\
%Affiliation \\
%Address \\
%\texttt{email} \\
%\And
%Coauthor \\
%Affiliation \\
%Address \\
%\texttt{email} \\
%\And
%Coauthor \\
%Affiliation \\
%Address \\
%\texttt{email} \\
%(if needed)\\
}
%\date{编辑于\today}\\
% The \author macro works with any number of authors. There are two commands
% used to separate the names and addresses of multiple authors: \And and \AND.
%
% Using \And between authors leaves it to \LaTeX{} to determine where to break
% the lines. Using \AND forces a linebreak at that point. So, if \LaTeX{}
% puts 3 of 4 authors names on the first line, and the last on the second
% line, try using \AND instead of \And before the third author name.

\newcommand{\fix}{\marginpar{FIX}}
\newcommand{\new}{\marginpar{NEW}}
\newcommand{\tabincell}[2]{\begin{tabular}{@{}#1@{}}#2\end{tabular}}
\nipsfinalcopy % Uncomment for camera-ready version
\newcommand\BackgroundPicture{%
  \put(0,0){%
    \parbox[b][\paperheight]{\paperwidth}{%
      \vfill
      \centering%
\begin{tikzpicture}[remember picture,overlay]
\node [rotate=45,scale=3,text opacity=0.2] at (current page.center) {CONFIDENTIAL GLASSSIX BENCHMARK};
\end{tikzpicture}%
      \vfill
    }}}
\begin{document}
\AddToShipoutPicture{\BackgroundPicture}

\maketitle

%\begin{abstract}
%    In order to evaluate the performance of face recognition, we established this benchmark to compare the internal released model, commercial solutions by our competitors and public models in academia.
%\end{abstract}
\section{Terminology Explanation(术语解释)}
\subsection{Approximate K-Nearest Neighbor(AKNN)}
AKNN算法通过计算不同特征值之间的距离，找出K个最相似(距离最近)的结果。
\subsection{Navigating Spread-out Graph(NSG)}
NSG算法属于AKNN的一种，测量距离找出K个最相似(距离最近)的结果。NSG的优点在于其具有局部性和单调性，这使得搜索速度得到很大提升，此算法已运用在淘宝的商品搜索中。
\subsection{Locality(局部性)}
邻居的邻居很可能也是邻居，根据这一想法，搜索最近邻的过程中，每次迭代仅在邻居的邻居范围内进行搜索，而不是整个数据集，这保证了搜索的速度。
\subsection{Monotonicity(单调性)}
NSG算法搜索过程具有单调性，每次搜索都距离目标更近一步，搜索路径没有环路，这保证了搜索的速度和准确度。
\subsection{Open Multi-Processing(OpenMP)}
OpenMP可以将程序由串行改为并行，充分利用CPU的多个计算核同时并行地进行计算，提升处理速度。例如在8核CPU上开启OpenMP后，计算速度理论上是不使用OpenMP相同程序的8倍。
\subsection{Single Instruction Multiple Data(SIMD)}
单指令多数据(SIMD)在收到一条指令后，可完成多条数据的处理任务，处理速度得以提升。不使用SIMD时，收到一条指令仅能处理一条数据。
\subsection{Build(构建)}
将大量的人脸特征值数据按照某种方式重新有序地组织，生成搜索结构(search structure)，以加快后续的搜索进程。
\subsection{Search(搜索)}
搜索是在搜索结构(search structure)上进行，返回K个最相似(距离最近)的结果,K值可自行修改。
\subsection{Accuracy(准确率)}
首先用耗时高的传统方法逐个比对人脸特征值数据，找出K个最相似(距离最近)的结果，假设为SetA；然后使用本算法在搜索结构(search structure)上找出K个最相似(距离最近)的结果，假设为SetB；最后计算SetB中结果出现在SetA的百分比，即为准确率。
\subsection{Maximum Memory Usage(峰值内存)}
程序运行过程中(包括构建和搜索)，监控内存占用情况,最大值即为峰值内存。
\subsection{128-Dimension Data Set(128维数据集)}
人脸特征值数据按128维格式存储，大量人脸特征值数据构成128维数据集。
\subsection{512-Dimension Data Set(512维数据集)}
人脸特征值数据按512维格式存储，大量人脸特征值数据构成512维数据集。
\section{Brief Introduction(算法简介)}
\subsection{Algorithm Description(算法描述)}
人脸搜索采用的NSG算法以图（graph）作为搜索结构(search structure)，并利用图（graph）的局部性和单调性来保证搜索的速度和准确度。本算法支持余弦距离和欧氏距离两种距离度量方式，并使用OpenMP和SIMD技术提升构建和搜索速度。
\subsection{Algorithm Processes(算法流程)}
\begin{enumerate}
  \item 根据人脸库数据使用AKNN算法，构建出过渡性搜索结构(search structure)——kgraph；
  \item 使用NSG算法在kgraph基础上再次进行构建，构建出具有单调性和局部性的搜索结构(search structure)——NSG，并保存在本地磁盘；
  \item 开始人脸搜索任务前，从本地磁盘加载NSG和人脸库数据；
  \item 对加载的NSG进行优化，将搜索结构(search structure)的数据对齐；
  \item 使用NSG开始人脸搜索，返回相似度最高的K个匹配结果。
\end{enumerate}
其中第1、2步合称为“构建”，第5步称为“搜索”。
\begin{figure}
\centering
\includegraphics[scale=0.4]{process.png}
\caption{NSG算法流程}
\label{1}
\end{figure}

\section{Benchmark Details(基准测试流程与测试数据)}
\subsection{Test Process(测试流程)}
测试主要在x86 平台和ARM 平台进行，测试NSG算法在不同平台、不同数据维度（128和512维）、不同数据量等条件下，构建时间、搜索时间、搜索准确率和峰值内存的情况。
然后测试Flann算法在相同条件下的性能表现，并与NSG算法进行对比分析。
\begin{enumerate}
  \item x86 平台（CPU: Intel E5-2640V4）测试；
  \item x86 平台（CPU: Intel i7-8700K）测试；
  \item x86 平台（CPU: Intel i5-8400）测试；
  \item ARM 平台（CPU: RK3399）测试。
\end{enumerate}
\subsection{Test Data(测试数据)}
测试数据集分为128维和512维的人脸特征值数据，512维数据由特征提取模块从人脸图片中生成；128维数据从512维数据中剪切得来，即512维数据的0-127位产生第一个128维数据，128-255位产生第二个128维数据，256-383位产生第三个128维数据，384-512位产生第四个128维数据。

根据实际使用场景，在x86 平台选取1K、2K、5K、10K、20K、50K、100K、200K、500K的数据量分别进行测试（K表示1000条），ARM平台测试数据量对应缩小10倍，分别为0.1K、0.2K、0.5K、1K、2K、5K、10K、20K、50K。
\section{x86 Platform Testing Protocols(x86 平台测试协议)}
\subsection{CPU: Intel E5-2640V4}
测试平台基本信息如表\ref{tab:basic-info-E5-2640V4}所示。
\begin{table}[H]
\caption{基本信息}
\label{tab:basic-info-E5-2640V4}
\begin{center}
\begin{tabular}{lc}
\multicolumn{1}{c}{\bf 细节}  &\multicolumn{1}{c}{\bf 描述}
\\ \hline \\
操作系统         &Windows 10 企业版 64位 \\
处理器           &Intel(R) Core(TM) E5-2640V4 CPU @ 2.40GHz \\
处理器核心数     &10核20线程 \\
内存             &32GB \\
编译器           &MSVC 19.5 \\
OpenMP           &开启 \\
SIMD指令集       &AVX2 \\
\end{tabular}
\end{center}
\end{table}


128维测试结果如表\ref{tab:E5-2640V4-NSG-128}所示。
\begin{table}[H]
\caption{$E5-2640V4-NSG-128$测试}
\label{tab:E5-2640V4-NSG-128}
\begin{center}
\begin{tabular}{ccccc}
\multicolumn{1}{c}{\bf 数据量}  &\multicolumn{1}{c}{\bf 构建时间(s)}  &\multicolumn{1}{c}{\bf 搜索时间(ms)}  &\multicolumn{1}{c}{\bf 准确率(\%)}  &\multicolumn{1}{c}{\bf 峰值内存(MB)}
\\ \hline \\
1k & 0 & 0 & 0 & 0\\
2k & 0 & 0 & 0 & 0\\
5k & 0 & 0 & 0 & 0\\
10k & 0 & 0 & 0 & 0\\
20k & 0 & 0 & 0 & 0\\
50k & 0 & 0 & 0 & 0\\
100k & 0 & 0 & 0 & 0\\
200k & 0 & 0 & 0 & 0\\
500k & 0 & 0 & 0 & 0\\
\end{tabular}
\end{center}
\end{table}

512维测试结果如表\ref{tab:E5-2640V4-NSG-512}所示。
\begin{table}[H]
\caption{$E5-2640V4-NSG-512$测试}
\label{tab:E5-2640V4-NSG-512}
\begin{center}
\begin{tabular}{ccccc}
\multicolumn{1}{c}{\bf 数据量}  &\multicolumn{1}{c}{\bf 构建时间(s)}  &\multicolumn{1}{c}{\bf 搜索时间(ms)}  &\multicolumn{1}{c}{\bf 准确率(\%)}  &\multicolumn{1}{c}{\bf 峰值内存(MB)}
\\ \hline \\
1k & 0 & 0 & 0 & 0\\
2k & 0 & 0 & 0 & 0\\
5k & 0 & 0 & 0 & 0\\
10k & 0 & 0 & 0 & 0\\
20k & 0 & 0 & 0 & 0\\
50k & 0 & 0 & 0 & 0\\
100k & 0 & 0 & 0 & 0\\
200k & 0 & 0 & 0 & 0\\
500k & 0 & 0 & 0 & 0\\
\end{tabular}
\end{center}
\end{table}

%\begin{table}[!htbp]
%\caption{$E5-2640V4-Flann-128$测试}
%\label{tab:E5-2640V4-Flann-128}
%\begin{center}
%\begin{tabular}{ccccc}
%\multicolumn{1}{c}{\bf 数据量}  &\multicolumn{1}{c}{\bf 构建时间(s)}  &\multicolumn{1}{c}{\bf 搜索时间(ms)}  &\multicolumn{1}{c}{\bf 准确率(\%)}  &\multicolumn{1}{c}{\bf 峰值内存(MB)}
%\\ \hline \\
%1k & 0 & 0 & 0 & 0\\
%2k & 0 & 0 & 0 & 0\\
%5k & 0 & 0 & 0 & 0\\
%10k & 0 & 0 & 0 & 0\\
%20k & 0 & 0 & 0 & 0\\
%50k & 0 & 0 & 0 & 0\\
%100k & 0 & 0 & 0 & 0\\
%200k & 0 & 0 & 0 & 0\\
%500k & 0 & 0 & 0 & 0\\
%\end{tabular}
%\end{center}
%\end{table}
%
%\begin{table}[!htbp]
%\caption{$E5-2640V4-Flann-512$测试}
%\label{tab:E5-2640V4-Flann-512}
%\begin{center}
%\begin{tabular}{ccccc}
%\multicolumn{1}{c}{\bf 数据量}  &\multicolumn{1}{c}{\bf 构建时间(s)}  &\multicolumn{1}{c}{\bf 搜索时间(ms)}  &\multicolumn{1}{c}{\bf 准确率(\%)}  &\multicolumn{1}{c}{\bf 峰值内存(MB)}
%\\ \hline \\
%1k & 0 & 0 & 0 & 0\\
%2k & 0 & 0 & 0 & 0\\
%5k & 0 & 0 & 0 & 0\\
%10k & 0 & 0 & 0 & 0\\
%20k & 0 & 0 & 0 & 0\\
%50k & 0 & 0 & 0 & 0\\
%100k & 0 & 0 & 0 & 0\\
%200k & 0 & 0 & 0 & 0\\
%500k & 0 & 0 & 0 & 0\\
%\end{tabular}
%\end{center}
%\end{table}

\subsection{CPU: Intel i7-8700K}
测试平台基本信息如表\ref{tab:basic-info-i7-8700k}所示。
\begin{table}[H]
\caption{基本信息}
\label{tab:basic-info-i7-8700k}
\begin{center}
\begin{tabular}{lc}
\multicolumn{1}{c}{\bf 细节}  &\multicolumn{1}{c}{\bf 描述}
\\ \hline \\
操作系统         &Windows 10 企业版 64位 \\
处理器           &Intel(R) Core(TM) i7-8700K CPU @ 3.70GHz \\
处理器核心数     &4核8线程 \\
内存             &32GB \\
编译器           &MSVC 19.5 \\
OpenMP           &开启 \\
SIMD指令集       &AVX2 \\
\end{tabular}
\end{center}
\end{table}

128维测试结果如表\ref{tab:i7-8700k-NSG-128}所示。
\begin{table}[H]
\caption{$i7-8700k-NSG-128$测试}
\label{tab:i7-8700k-NSG-128}
\begin{center}
\begin{tabular}{ccccc}
\multicolumn{1}{c}{\bf 数据量}  &\multicolumn{1}{c}{\bf 构建时间(s)}  &\multicolumn{1}{c}{\bf 搜索时间(ms)}  &\multicolumn{1}{c}{\bf 准确率(\%)}  &\multicolumn{1}{c}{\bf 峰值内存(MB)}
\\ \hline \\
1k & 0.12 & 0.06 & 100 & 6\\
2k & 0.13 & 0.08 & 100 & 11\\
5k & 0.55 & 0.10 & 99.94 & 22\\
10k & 1.12 & 0.11 & 100 & 41\\
20k & 2.05 & 0.13 & 99.97 & 68\\
50k & 5.48 & 0.15 & 99.52 & 177\\
100k & 16.05 & 0.17 & 99.84 & 375\\
200k & 54.29 & 0.25 & 99.94 & 753\\
500k & 227.5 & 0.36 & 98.88 & 2008\\
\end{tabular}
\end{center}
\end{table}

512维测试结果如表\ref{tab:i7-8700k-NSG-512}所示。
\begin{table}[H]
\caption{$i7-8700k-NSG-512$测试}
\label{tab:i7-8700k-NSG-512}
\begin{center}
\begin{tabular}{ccccc}
\multicolumn{1}{c}{\bf 数据量}  &\multicolumn{1}{c}{\bf 构建时间(s)}  &\multicolumn{1}{c}{\bf 搜索时间(ms)}  &\multicolumn{1}{c}{\bf 准确率(\%)}  &\multicolumn{1}{c}{\bf 峰值内存(MB)}
\\ \hline \\
1k & 0.23 & 0.13 & 100 & 13\\
2k & 0.42 & 0.14 & 100 & 23\\
5k & 1.67 & 0.18 & 100 & 55\\
10k & 3.22 & 0.21 & 99.93 & 86\\
20k & 6.82 & 0.27 & 99.95 & 174\\
50k & 18.95 & 0.30 & 99.94 & 499\\
100k & 59.7 & 0.41 & 99.80 & 820\\
200k & 164.3 & 0.56 & 99.75 & 1674\\
500k & 588.3 & 0.63 & 99.15 & 3597\\
\end{tabular}
\end{center}
\end{table}

%\begin{table}[!htbp]
%\caption{$i7-8700k-Flann-128$测试}
%\label{tab:i7-8700k-Flann-128}
%\begin{center}
%\begin{tabular}{ccccc}
%\multicolumn{1}{c}{\bf 数据量}  &\multicolumn{1}{c}{\bf 构建时间(s)}  &\multicolumn{1}{c}{\bf 搜索时间(ms)}  &\multicolumn{1}{c}{\bf 准确率(\%)}  &\multicolumn{1}{c}{\bf 峰值内存(MB)}
%\\ \hline \\
%1k & 0 & 0 & 0 & 0\\
%2k & 0 & 0 & 0 & 0\\
%5k & 0 & 0 & 0 & 0\\
%10k & 0 & 0 & 0 & 0\\
%20k & 0 & 0 & 0 & 0\\
%50k & 0 & 0 & 0 & 0\\
%100k & 0 & 0 & 0 & 0\\
%200k & 0 & 0 & 0 & 0\\
%500k & 0 & 0 & 0 & 0\\
%\end{tabular}
%\end{center}
%\end{table}
%
%\begin{table}[!htbp]
%\caption{$i7-8700k-Flann-512$测试}
%\label{tab:i7-8700k-Flann-512}
%\begin{center}
%\begin{tabular}{ccccc}
%\multicolumn{1}{c}{\bf 数据量}  &\multicolumn{1}{c}{\bf 构建时间(s)}  &\multicolumn{1}{c}{\bf 搜索时间(ms)}  &\multicolumn{1}{c}{\bf 准确率(\%)}  &\multicolumn{1}{c}{\bf 峰值内存(MB)}
%\\ \hline \\
%1k & 0 & 0 & 0 & 0\\
%2k & 0 & 0 & 0 & 0\\
%5k & 0 & 0 & 0 & 0\\
%10k & 0 & 0 & 0 & 0\\
%20k & 0 & 0 & 0 & 0\\
%50k & 0 & 0 & 0 & 0\\
%100k & 0 & 0 & 0 & 0\\
%200k & 0 & 0 & 0 & 0\\
%500k & 0 & 0 & 0 & 0\\
%\end{tabular}
%\end{center}
%\end{table}

\subsection{CPU: Intel i5-8400}
测试平台基本信息如表\ref{tab:basic-info-i5-8400}所示。
\begin{table}[H]
\caption{基本信息}
\label{tab:basic-info-i5-8400}
\begin{center}
\begin{tabular}{lc}
\multicolumn{1}{c}{\bf 细节}  &\multicolumn{1}{c}{\bf 描述}
\\ \hline \\
操作系统         &Windows 10 企业版 64位 \\
处理器           &Intel(R) Core(TM) i5-8400 CPU @ 2.80GHz \\
处理器核心数     &4核4线程 \\
内存             &32GB \\
编译器           &MSVC 19.5 \\
OpenMP           &开启 \\
SIMD指令集       &AVX2 \\
\end{tabular}
\end{center}
\end{table}

128维测试结果如表\ref{tab:i5-8400-NSG-128}所示。
\begin{table}[H]
\caption{$i5-8400-NSG-128$测试}
\label{tab:i5-8400-NSG-128}
\begin{center}
\begin{tabular}{ccccc}
\multicolumn{1}{c}{\bf 数据量}  &\multicolumn{1}{c}{\bf 构建时间(s)}  &\multicolumn{1}{c}{\bf 搜索时间(ms)}  &\multicolumn{1}{c}{\bf 准确率(\%)}  &\multicolumn{1}{c}{\bf 峰值内存(MB)}
\\ \hline \\
1k & 0.16 & 0.07 & 100 & 6\\
2k & 0.22 & 0.10 & 100 & 10\\
5k & 1.00 & 0.12 & 99.92 & 20\\
10k & 1.89 & 0.13 & 100 & 39\\
20k & 3.21 & 0.14 & 99.95 & 66\\
50k & 8.55 & 0.16 & 99.63 & 174\\
100k & 25.04 & 0.18 & 99.85 & 373\\
200k & 81.15 & 0.27 & 99.95 & 749\\
500k & 362.32 & 0.39 & 98.88 & 2002\\
\end{tabular}
\end{center}
\end{table}

512维测试结果如表\ref{tab:i5-8400-NSG-512}所示。
\begin{table}[H]
\caption{$i5-8400-NSG-512$测试}
\label{tab:i5-8400-NSG-512}
\begin{center}
\begin{tabular}{ccccc}
\multicolumn{1}{c}{\bf 数据量}  &\multicolumn{1}{c}{\bf 构建时间(s)}  &\multicolumn{1}{c}{\bf 搜索时间(ms)}  &\multicolumn{1}{c}{\bf 准确率(\%)}  &\multicolumn{1}{c}{\bf 峰值内存(MB)}
\\ \hline \\
1k & 0.37 & 0.12 & 100 & 11\\
2k & 0.73 & 0.14 & 100 & 21\\
5k & 2.59 & 0.17 & 100 & 53\\
10k & 5.14 & 0.21 & 99.93 & 83\\
20k & 10.77 & 0.26 & 99.95 & 171\\
50k & 29.40 & 0.32 & 99.95 & 494\\
100k & 87.76 & 0.43 & 99.81 & 816\\
200k & 233.75 & 0.61 & 99.75 & 1670\\
500k & 845.38 & 0.68 & 99.14 & 3595\\
\end{tabular}
\end{center}
\end{table}

%\begin{table}[!htbp]
%\caption{$i5-8400-Flann-128$测试}
%\label{tab:i5-8400-Flann-128}
%\begin{center}
%\begin{tabular}{ccccc}
%\multicolumn{1}{c}{\bf 数据量}  &\multicolumn{1}{c}{\bf 构建时间(s)}  &\multicolumn{1}{c}{\bf 搜索时间(ms)}  &\multicolumn{1}{c}{\bf 准确率(\%)}  &\multicolumn{1}{c}{\bf 峰值内存(MB)}
%\\ \hline \\
%1k & 0 & 0 & 0 & 0\\
%2k & 0 & 0 & 0 & 0\\
%5k & 0 & 0 & 0 & 0\\
%10k & 0 & 0 & 0 & 0\\
%20k & 0 & 0 & 0 & 0\\
%50k & 0 & 0 & 0 & 0\\
%100k & 0 & 0 & 0 & 0\\
%200k & 0 & 0 & 0 & 0\\
%500k & 0 & 0 & 0 & 0\\
%\end{tabular}
%\end{center}
%\end{table}
%
%\begin{table}[!htbp]
%\caption{$i5-8400-Flann-512$测试}
%\label{tab:i5-8400-Flann-512}
%\begin{center}
%\begin{tabular}{ccccc}
%\multicolumn{1}{c}{\bf 数据量}  &\multicolumn{1}{c}{\bf 构建时间(s)}  &\multicolumn{1}{c}{\bf 搜索时间(ms)}  &\multicolumn{1}{c}{\bf 准确率(\%)}  &\multicolumn{1}{c}{\bf 峰值内存(MB)}
%\\ \hline \\
%1k & 0 & 0 & 0 & 0\\
%2k & 0 & 0 & 0 & 0\\
%5k & 0 & 0 & 0 & 0\\
%10k & 0 & 0 & 0 & 0\\
%20k & 0 & 0 & 0 & 0\\
%50k & 0 & 0 & 0 & 0\\
%100k & 0 & 0 & 0 & 0\\
%200k & 0 & 0 & 0 & 0\\
%500k & 0 & 0 & 0 & 0\\
%\end{tabular}
%\end{center}
%\end{table}

\section{ARM Platform Testing Protocols(ARM平台测试协议)}
\subsection{CPU: RK3399}

\section{Legacy Issues(遗留问题)}

\begin{enumerate}
  \item 构建时间不能满足实时性要求；
  \item 研究是否能使用CUDA完成部分distance计算任务。
\end{enumerate}
%Papers to be submitted to NIPS 2013 must be prepared according to the
%instructions presented here. Papers may be only up to eight pages long,
%including figures. Since 2009 an additional ninth page \textit{containing only
%cited references} is allowed. Papers that exceed nine pages will not be
%reviewed, or in any other way considered for presentation at the conference.
%%This is a strict upper bound.
%
%Please note that this year we have introduced automatic line number generation
%into the style file (for \LaTeXe and Word versions). This is to help reviewers
%refer to specific lines of the paper when they make their comments. Please do
%NOT refer to these line numbers in your paper as they will be removed from the
%style file for the final version of accepted papers.
%
%The margins in 2013 are the same as since 2007, which allow for $\approx 15\%$
%more words in the paper compared to earlier years. We are also again using
%double-blind reviewing. Both of these require the use of new style files.
%
%Authors are required to use the NIPS \LaTeX{} style files obtainable at the
%NIPS website as indicated below. Please make sure you use the current files and
%not previous versions. Tweaking the style files may be grounds for rejection.
%
%%% \subsection{Double-blind reviewing}
%
%%% This year we are doing double-blind reviewing: the reviewers will not know
%%% who the authors of the paper are. For submission, the NIPS style file will
%%% automatically anonymize the author list at the beginning of the paper.
%
%%% Please write your paper in such a way to preserve anonymity. Refer to
%%% previous work by the author(s) in the third person, rather than first
%%% person. Do not provide Web links to supporting material at an identifiable
%%% web site.
%
%%%\subsection{Electronic submission}
%%%
%%% \textbf{THE SUBMISSION DEADLINE IS MAY 31st, 2013. SUBMISSIONS MUST BE LOGGED BY
%%% 23:00, MAY 31st, 2013, UNIVERSAL TIME}
%
%%% You must enter your submission in the electronic submission form available at
%%% the NIPS website listed above. You will be asked to enter paper title, name of
%%% all authors, keyword(s), and data about the contact
%%% author (name, full address, telephone, fax, and email). You will need to
%%% upload an electronic (postscript or pdf) version of your paper.
%
%%% You can upload more than one version of your paper, until the
%%% submission deadline. We strongly recommended uploading your paper in
%%% advance of the deadline, so you can avoid last-minute server congestion.
%%%
%%% Note that your submission is only valid if you get an e-mail
%%% confirmation from the server. If you do not get such an e-mail, please
%%% try uploading again.
%
%
%\subsection{Retrieval of style files}
%
%The style files for NIPS and other conference information are available on the World Wide Web at
%\begin{center}
%   \url{http://www.nips.cc/}
%\end{center}
%The file \verb+nips2013.pdf+ contains these
%instructions and illustrates the
%various formatting requirements your NIPS paper must satisfy. \LaTeX{}
%users can choose between two style files:
%\verb+nips11submit_09.sty+ (to be used with \LaTeX{} version 2.09) and
%\verb+nips11submit_e.sty+ (to be used with \LaTeX{}2e). The file
%\verb+nips2013.tex+ may be used as a ``shell'' for writing your paper. All you
%have to do is replace the author, title, abstract, and text of the paper with
%your own. The file
%\verb+nips2013.rtf+ is provided as a shell for MS Word users.
%
%The formatting instructions contained in these style files are summarized in
%sections \ref{gen_inst}, \ref{headings}, and \ref{others} below.
%
%%% \subsection{Keywords for paper submission}
%%% Your NIPS paper can be submitted with any of the following keywords (more than one keyword is possible for each paper):
%
%%% \begin{verbatim}
%%% Bioinformatics
%%% Biological Vision
%%% Brain Imaging and Brain Computer Interfacing
%%% Clustering
%%% Cognitive Science
%%% Control and Reinforcement Learning
%%% Dimensionality Reduction and Manifolds
%%% Feature Selection
%%% Gaussian Processes
%%% Graphical Models
%%% Hardware Technologies
%%% Kernels
%%% Learning Theory
%%% Machine Vision
%%% Margins and Boosting
%%% Neural Networks
%%% Neuroscience
%%% Other Algorithms and Architectures
%%% Other Applications
%%% Semi-supervised Learning
%%% Speech and Signal Processing
%%% Text and Language Applications
%
%%% \end{verbatim}
%
%\section{General formatting instructions}
%\label{gen_inst}
%
%The text must be confined within a rectangle 5.5~inches (33~picas) wide and
%9~inches (54~picas) long. The left margin is 1.5~inch (9~picas).
%Use 10~point type with a vertical spacing of 11~points. Times New Roman is the
%preferred typeface throughout. Paragraphs are separated by 1/2~line space,
%with no indentation.
%
%Paper title is 17~point, initial caps/lower case, bold, centered between
%2~horizontal rules. Top rule is 4~points thick and bottom rule is 1~point
%thick. Allow 1/4~inch space above and below title to rules. All pages should
%start at 1~inch (6~picas) from the top of the page.
%
%%The version of the paper submitted for review should have ``Anonymous Author(s)'' as the author of the paper.
%
%For the final version, authors' names are
%set in boldface, and each name is centered above the corresponding
%address. The lead author's name is to be listed first (left-most), and
%the co-authors' names (if different address) are set to follow. If
%there is only one co-author, list both author and co-author side by side.
%
%Please pay special attention to the instructions in section \ref{others}
%regarding figures, tables, acknowledgments, and references.
%
%\section{Headings: first level}
%\label{headings}
%
%First level headings are lower case (except for first word and proper nouns),
%flush left, bold and in point size 12. One line space before the first level
%heading and 1/2~line space after the first level heading.
%
%\subsection{Headings: second level}
%
%Second level headings are lower case (except for first word and proper nouns),
%flush left, bold and in point size 10. One line space before the second level
%heading and 1/2~line space after the second level heading.
%
%\subsubsection{Headings: third level}
%
%Third level headings are lower case (except for first word and proper nouns),
%flush left, bold and in point size 10. One line space before the third level
%heading and 1/2~line space after the third level heading.
%
%\section{Citations, figures, tables, references}
%\label{others}
%
%These instructions apply to everyone, regardless of the formatter being used.
%
%\subsection{Citations within the text}
%
%Citations within the text should be numbered consecutively. The corresponding
%number is to appear enclosed in square brackets, such as [1] or [2]-[5]. The
%corresponding references are to be listed in the same order at the end of the
%paper, in the \textbf{References} section. (Note: the standard
%\textsc{Bib\TeX} style \texttt{unsrt} produces this.) As to the format of the
%references themselves, any style is acceptable as long as it is used
%consistently.
%
%As submission is double blind, refer to your own published work in the
%third person. That is, use ``In the previous work of Jones et al.\ [4]'',
%not ``In our previous work [4]''. If you cite your other papers that
%are not widely available (e.g.\ a journal paper under review), use
%anonymous author names in the citation, e.g.\ an author of the
%form ``A.\ Anonymous''.
%
%
%\subsection{Footnotes}
%
%Indicate footnotes with a number\footnote{Sample of the first footnote} in the
%text. Place the footnotes at the bottom of the page on which they appear.
%Precede the footnote with a horizontal rule of 2~inches
%(12~picas).\footnote{Sample of the second footnote}
%
%\subsection{Figures}
%
%All artwork must be neat, clean, and legible. Lines should be dark
%enough for purposes of reproduction; art work should not be
%hand-drawn. The figure number and caption always appear after the
%figure. Place one line space before the figure caption, and one line
%space after the figure. The figure caption is lower case (except for
%first word and proper nouns); figures are numbered consecutively.
%
%Make sure the figure caption does not get separated from the figure.
%Leave sufficient space to avoid splitting the figure and figure caption.
%
%You may use color figures.
%However, it is best for the
%figure captions and the paper body to make sense if the paper is printed
%either in black/white or in color.
%\begin{figure}[h]
%\begin{center}
%%\framebox[4.0in]{$\;$}
%%\fbox{\rule[-.5cm]{0cm}{4cm} \rule[-.5cm]{4cm}{0cm}}
%\includegraphics[width=2in]{logo.png}
%\end{center}
%\caption{Sample figure caption.}
%\end{figure}
%
%\subsection{Tables}
%
%All tables must be centered, neat, clean and legible. Do not use hand-drawn
%tables. The table number and title always appear before the table. See
%Table~\ref{sample-table}.
%
%Place one line space before the table title, one line space after the table
%title, and one line space after the table. The table title must be lower case
%(except for first word and proper nouns); tables are numbered consecutively.
%
%\begin{table}[t]
%\caption{Sample table title}
%\label{sample-table}
%\begin{center}
%\begin{tabular}{ll}
%\multicolumn{1}{c}{\bf PART}  &\multicolumn{1}{c}{\bf DESCRIPTION}
%\\ \hline \\
%Dendrite         &Input terminal \\
%Axon             &Output terminal \\
%Soma             &Cell body (contains cell nucleus) \\
%\end{tabular}
%\end{center}
%\end{table}
%
%\section{Final instructions}
%Do not change any aspects of the formatting parameters in the style files.
%In particular, do not modify the width or length of the rectangle the text
%should fit into, and do not change font sizes (except perhaps in the
%\textbf{References} section; see below). Please note that pages should be
%numbered.
%
%\section{Preparing PostScript or PDF files}
%
%Please prepare PostScript or PDF files with paper size ``US Letter'', and
%not, for example, ``A4''. The -t
%letter option on dvips will produce US Letter files.
%
%Fonts were the main cause of problems in the past years. Your PDF file must
%only contain Type 1 or Embedded TrueType fonts. Here are a few instructions
%to achieve this.
%
%\begin{itemize}
%
%\item You can check which fonts a PDF files uses.  In Acrobat Reader,
%select the menu Files$>$Document Properties$>$Fonts and select Show All Fonts. You can
%also use the program \verb+pdffonts+ which comes with \verb+xpdf+ and is
%available out-of-the-box on most Linux machines.
%
%\item The IEEE has recommendations for generating PDF files whose fonts
%are also acceptable for NIPS. Please see
%\url{http://www.emfield.org/icuwb2010/downloads/IEEE-PDF-SpecV32.pdf}
%
%\item LaTeX users:
%
%\begin{itemize}
%
%\item Consider directly generating PDF files using \verb+pdflatex+
%(especially if you are a MiKTeX user).
%PDF figures must be substituted for EPS figures, however.
%
%\item Otherwise, please generate your PostScript and PDF files with the following commands:
%\begin{verbatim}
%dvips mypaper.dvi -t letter -Ppdf -G0 -o mypaper.ps
%ps2pdf mypaper.ps mypaper.pdf
%\end{verbatim}
%
%Check that the PDF files only contains Type 1 fonts.
%%For the final version, please send us both the Postscript file and
%%the PDF file.
%
%\item xfig "patterned" shapes are implemented with
%bitmap fonts.  Use "solid" shapes instead.
%\item The \verb+\bbold+ package almost always uses bitmap
%fonts.  You can try the equivalent AMS Fonts with command
%\begin{verbatim}
%\usepackage[psamsfonts]{amssymb}
%\end{verbatim}
% or use the following workaround for reals, natural and complex:
%\begin{verbatim}
%\newcommand{\RR}{I\!\!R} %real numbers
%\newcommand{\Nat}{I\!\!N} %natural numbers
%\newcommand{\CC}{I\!\!\!\!C} %complex numbers
%\end{verbatim}
%
%\item Sometimes the problematic fonts are used in figures
%included in LaTeX files. The ghostscript program \verb+eps2eps+ is the simplest
%way to clean such figures. For black and white figures, slightly better
%results can be achieved with program \verb+potrace+.
%\end{itemize}
%\item MSWord and Windows users (via PDF file):
%\begin{itemize}
%\item Install the Microsoft Save as PDF Office 2007 Add-in from
%\url{http://www.microsoft.com/downloads/details.aspx?displaylang=en\&familyid=4d951911-3e7e-4ae6-b059-a2e79ed87041}
%\item Select ``Save or Publish to PDF'' from the Office or File menu
%\end{itemize}
%\item MSWord and Mac OS X users (via PDF file):
%\begin{itemize}
%\item From the print menu, click the PDF drop-down box, and select ``Save
%as PDF...''
%\end{itemize}
%\item MSWord and Windows users (via PS file):
%\begin{itemize}
%\item To create a new printer
%on your computer, install the AdobePS printer driver and the Adobe Distiller PPD file from
%\url{http://www.adobe.com/support/downloads/detail.jsp?ftpID=204} {\it Note:} You must reboot your PC after installing the
%AdobePS driver for it to take effect.
%\item To produce the ps file, select ``Print'' from the MS app, choose
%the installed AdobePS printer, click on ``Properties'', click on ``Advanced.''
%\item Set ``TrueType Font'' to be ``Download as Softfont''
%\item Open the ``PostScript Options'' folder
%\item Select ``PostScript Output Option'' to be ``Optimize for Portability''
%\item Select ``TrueType Font Download Option'' to be ``Outline''
%\item Select ``Send PostScript Error Handler'' to be ``No''
%\item Click ``OK'' three times, print your file.
%\item Now, use Adobe Acrobat Distiller or ps2pdf to create a PDF file from
%the PS file. In Acrobat, check the option ``Embed all fonts'' if
%applicable.
%\end{itemize}
%
%\end{itemize}
%If your file contains Type 3 fonts or non embedded TrueType fonts, we will
%ask you to fix it.
%
%\subsection{Margins in LaTeX}
%
%Most of the margin problems come from figures positioned by hand using
%\verb+\special+ or other commands. We suggest using the command
%\verb+\includegraphics+
%from the graphicx package. Always specify the figure width as a multiple of
%the line width as in the example below using .eps graphics
%\begin{verbatim}
%   \usepackage[dvips]{graphicx} ...
%   \includegraphics[width=0.8\linewidth]{myfile.eps}
%\end{verbatim}
%or % Apr 2009 addition
%\begin{verbatim}
%   \usepackage[pdftex]{graphicx} ...
%   \includegraphics[width=0.8\linewidth]{myfile.pdf}
%\end{verbatim}
%for .pdf graphics.
%See section 4.4 in the graphics bundle documentation (\url{http://www.ctan.org/tex-archive/macros/latex/required/graphics/grfguide.ps})
%
%A number of width problems arise when LaTeX cannot properly hyphenate a
%line. Please give LaTeX hyphenation hints using the \verb+\-+ command.
%
%
%\subsubsection*{Acknowledgments}
%
%Use unnumbered third level headings for the acknowledgments. All
%acknowledgments go at the end of the paper. Do not include
%acknowledgments in the anonymized submission, only in the
%final paper.
%
%\subsubsection*{References}
%
%References follow the acknowledgments. Use unnumbered third level heading for
%the references. Any choice of citation style is acceptable as long as you are
%consistent. It is permissible to reduce the font size to `small' (9-point)
%when listing the references. {\bf Remember that this year you can use
%a ninth page as long as it contains \emph{only} cited references.}
%
%\small{
%[1] Alexander, J.A. \& Mozer, M.C. (1995) Template-based algorithms
%for connectionist rule extraction. In G. Tesauro, D. S. Touretzky
%and T.K. Leen (eds.), {\it Advances in Neural Information Processing
%Systems 7}, pp. 609-616. Cambridge, MA: MIT Press.
%
%[2] Bower, J.M. \& Beeman, D. (1995) {\it The Book of GENESIS: Exploring
%Realistic Neural Models with the GEneral NEural SImulation System.}
%New York: TELOS/Springer-Verlag.
%
%[3] Hasselmo, M.E., Schnell, E. \& Barkai, E. (1995) Dynamics of learning
%and recall at excitatory recurrent synapses and cholinergic modulation
%in rat hippocampal region CA3. {\it Journal of Neuroscience}
%{\bf 15}(7):5249-5262.
%}

\end{document}
